\documentclass[12pt]{article}
\usepackage[utf8]{inputenc}
\usepackage[T1]{fontenc}
\usepackage[margin=1in]{geometry}
\usepackage{setspace}
\usepackage{xcolor}
\usepackage[colorlinks=true, urlcolor=blue!60!black, linkcolor=blue!60!black]{hyperref}
\setlength{\parindent}{0em}
\setlength{\parskip}{0.45em}
\renewcommand{\baselinestretch}{1.0}
\begin{document}
\begin{center}
{\Large \textit{United States v. Texas}, 599 U.S. 670 (2023)\par}
\vspace{0.3em}
{\normalsize Case Brief --- One-Page Edition\par}
\end{center}

\textbf{Memory Jogger.} States lack standing to make the Executive arrest more noncitizens; indirect costs of non-enforcement are not judicially cognizable (\href{https://www.law.cornell.edu/supremecourt/text/22-58}{Cornell LII}).

\textbf{Facts.} DHS's 2021 Guidelines prioritized noncitizens who threaten national security, public safety, or border security. Texas and Louisiana said the Guidelines violated statutes providing that DHS ``shall'' take custody of certain criminal noncitizens, costing the States money for supervision, recidivism, education, and healthcare.

\textbf{Procedural History.} The district court vacated the Guidelines nationwide after a bench trial; the Court granted certiorari before judgment.

\textbf{Votes.} Reversed, 8--1: Kavanaugh for five; Gorsuch (with Thomas and Barrett) and Barrett concurred in the judgment; Alito dissented.

\textbf{Holding.} A State neither prosecuted nor threatened with prosecution cannot challenge arrest and prosecution policies, even if it bears indirect costs. Legality was not decided.

\textbf{Reasoning.} No tradition supports orders compelling arrests (\textit{Linda R. S.}); non-enforcement coerces no one, intrudes on Article II discretion, and turns on resource trade-offs; five possible exceptions were reserved. Gorsuch found no redressability and doubted universal vacatur; Barrett faulted the majority's reading of precedent. Alito said Texas proved real costs and the ruling revives royal suspending powers.

\textbf{Examples.} \emph{Same-Side:} a state contesting deprioritized marijuana prosecutions. \emph{Opposite-Side:} a program granting work authorization. \emph{Fence-Sitter:} a statute authorizing state suits to compel custody.

\textbf{Critique.} It curbs state attorney-general suits but, critics say, lets agencies disregard mandatory statutes.

\textbf{Key Quotation.} ``The States essentially want the Federal Judiciary to order the Executive Branch to alter its arrest policy so as to make more arrests. But this Court has long held `that a citizen lacks standing to contest the policies of the prosecuting authority when he himself is neither prosecuted nor threatened with prosecution.' \textit{Linda R. S. v. Richard D.}, 410 U. S. 614, 619 (1973). Consistent with that fundamental Article III principle, we conclude that the States lack Article III standing to bring this suit.'' --- Justice Kavanaugh, for the Court.

\end{document}
